This section explains a couple of useful things about Markov chains. When not otherwise specified Markov chains are assumed to be of first-order. 

Let us start with a definition: $q_t$ is a (discrete-time) $n$-state first-order Markov chain if at any time $t$, $q_t$ takes a value in $\{q_1,...,q_n\}$, and
the probability that $q$ takes the value $q_j$ at time $t+1$ depends only on the state of $q$ at time $t$, that is $Pr(q_{t+1}=q_j)=Pr(q_{t+1}=q_i | q_t=q_j)$.

The transition matrix of $q$ is defined as the $n$-by-$n$ matrix $\Pi^q=[\pi^q_{ij}]$, where $\pi^q_{ij}=Pr(q_{t+1}=q_i | q_t=p_j)$. Thus the element of the
transition matrix in row $i$ and column $j$ gives the probability that the state tomorrow is $q_j$ given that the state today is $q_i$.

For a more comprehensive treatment of Markov Chains see SLP chapters 8, 11, \& 12; Sargent \& Ljungquist chapter 2; or Grimmett \& Stirzaker - Probability
and Random Processes.

\subsection{How to turn a Markov chain with two variables into a single variable Markov chain}
Say you have a Markov chain with two variables, $q$ \& $p$ (the most common macroeconomic application for this is in heterogenous models, where $q$ is an 
idiosyncratic state $s$ and $p$ is the aggregate state $z$). Let $q$ take the states $q_1,...q_n$, and $p$ take the states $p_1,...,p_m$. We start with 
the simple example of how to combine the two when their transitions are completely independent of each other so as to illustrate the concepts and then 
treat the general case.

When $q$ and $p$ are independent Markov chains, $Pr(q_{t+1}=q_j)=Pr(q_{t+1}=q_i | q_t=q_j)=\pi^q_{ij}$ $\forall i,j=1,...n$ and 
$Pr(p_{t+1}=p_j)=Pr(p_{t+1}=p_i | p_t=p_j)=\pi^p_{ij}$ $\forall i,j=1,...m$, then we can define a new single-variable Markov chain $r$
simply by taking the Kronecker Product of $q$ and $p$. Thus, $r$ will have $n$ times $m$ states, $[r_1,..,r_{nm}]'=[q_1,...q_n]'\otimes[p_1,...p_m]'$, and it's transition matrix
will be the Kronecker Product of their transition matrices; $\Pi^r=\Pi^q \otimes \Pi^p$. For the definition of the Kronecker product of two matrices see
wikipedia.

For the vector $(q,p)$ to be a (first-order) Markov chain, at least one of $q$ and $p$ must be independent of the current period value of the other. Thus 
we assume, without loss of generality, that $q_t$ evolves according to the transformation matrix defined by $Pr(q_t)=Pr(q_t | q_{t-1},p_{t-1})$, while $p_t$ 
has the transformation matrix $Pr(p_t)=Pr(q_t, q_{t-1}, p_{t-1})$\footnote{Note that here I switch from describing Markov chains as $t+1 |t$ to $t|t-1$, the
difference is purely cosmetic.}. Again we can define a Markov chain $r$ will have $n$ times $m$ states by $[r_1,..,r_{nm}]'=[q_1,...q_n]'\otimes[p_1,...p_m]'$.
The transition matrix however is now more complicated; rather than provide a general formula you are referred to the examples of Imrohoroglu (1989) and
OTHEREXAMPLE which illustrate some of the common cases you may wish to model.

If you have three or more variables in vector of the first-order Markov chain you can reduce this to a scalar first-order Markov chain simply by iteratively
combining pairs. For example with three variables, $(q^1, q^2, q^3)$, start by first defining $r^1$ as the combination of $q^1$ \& $q^2$ (combining them as
described above), and then $r$ as the combination of $r^1$ and $q^3$.


\subsection{How to turn a second-order Markov chain into a first-order Markov chain}
Suppose that $q$, with states $q_1,...q_n$ is a second-order Markov chain, that is $Pr(q_{t+1}=q_i)=Pr(q_{t+1}=q_i|q_t=q_j, q_{t-1}=q_k)
\neq Pr(q_{t+1}=q_i|q_t=q_j)$. Consider now the vector $(q_{t+1}, q_t)$. It turns out that this vector is in fact a vector first-order Markov
chain (define this periods state $(q_{t+1}, q_t)$ in terms of last periods state $(q_t,q_{t-1})$ by defining this periods $q_{t+1}$ as a function of
last periods $q_t$ \& $q_{t-1}$ following the original second-order Markov chain, and define this periods $q_t$ as being last periods $q_t$). Now just
combine the two elements of this first-order vector Markov chain, $(q_{t+1}, q_t)$, into a scalar first-order Markov chain $r$ in the same way as we
did above, that is by the Kronecker product.

For third- and higher-order Markov chains simply turn them into three- and higher-element first-order vector Markov chains, and then combine them 
into scalars by repeated pairwise Kronecker products as above. 
